\documentclass{article}
\usepackage{amsmath,amssymb}
\usepackage{xurl}
\usepackage[hidelinks]{hyperref}
% Shared commands for notes in docs/paper-gaps/.

\DeclareUnicodeCharacter{2080}{\ifmmode{}_{0}\else\textsubscript{0}\fi}
\DeclareUnicodeCharacter{2081}{\ifmmode{}_{1}\else\textsubscript{1}\fi}
\DeclareUnicodeCharacter{2082}{\ifmmode{}_{2}\else\textsubscript{2}\fi}
\DeclareUnicodeCharacter{2099}{\ifmmode{}_{n}\else\textsubscript{n}\fi}

\newcommand{\Lean}{\textsc{Lean}}
\ExplSyntaxOn
\NewDocumentCommand{\leanid}{m}{
  \ifmmode
    \textsf{\detokenize{#1}}
  \else
    \group_begin:
    \urlstyle{sf}
    \tl_set:Nn \l_tmpa_tl {#1}
    \tl_replace_all:Nnn \l_tmpa_tl {\_} {_}
    \exp_args:NV \path \l_tmpa_tl
    \group_end:
  \fi
}
\ExplSyntaxOff
\providecommand{\tr}{\operatorname{tr}}
\newcommand{\tnleanIssueBase}{https://github.com/LionSR/TNLean/issues/}
\newcommand{\tnleanPrBase}{https://github.com/LionSR/TNLean/pull/}
\newcommand{\tnleanDocBase}{https://sirui-lu.com/TNLean/docs/}
\newcommand{\ghissue}[1]{\href{\tnleanIssueBase#1}{GitHub issue \##1}}
\newcommand{\ghpr}[1]{\href{\tnleanPrBase#1}{PR \##1}}
\newcommand{\doclink}[1]{\href{\tnleanDocBase#1.html}{\path{#1}}}

% Dirac notation and the identity operator, as in blueprint/src/macros/common.tex.
% \providecommand keeps note-local definitions authoritative.
\usepackage{dsfont}
\providecommand{\ket}[1]{|#1\rangle}
\providecommand{\bra}[1]{\langle #1|}
\providecommand{\braket}[2]{\langle #1 | #2 \rangle}
\providecommand{\ketbra}[2]{|#1\rangle\!\langle #2|}
\providecommand{\Id}{\mathds{1}}
\providecommand{\one}{\mathds{1}}
\providecommand{\C}{\mathbb{C}}
\providecommand{\MN}[1]{M_{#1}(\C)}
\providecommand{\MD}{M_D(\C)}

% Verdict marker: \gapnote{<kind>}{<status>}. Kind is the policy
% classification (clarification, local-correction, scope-restriction,
% unfaithful, false-source, open-gap); status is open, wip (elimination
% underway), resolved, or historical. Machine-read by the site index;
% prints a verdict line.
\newcommand{\gapnote}[2]{\par\noindent\textbf{Verdict:} #1 (#2).\par}

\title{Quantum-double Hamiltonian terms: normalization and local scope}
\author{}
\date{2026-10-05}
\begin{document}
\maketitle
\gapnote{local-correction}{resolved}

\section{Source and normalization correction}
Schuch, Cirac, and P\'erez-Garc\'ia,
\href{https://arxiv.org/abs/1001.3807}{arXiv:1001.3807v3},
\path{Papers/1001.3807/paper_v3.tex}, lines 2861--2877, gives the
product-one diagonal plaquette term and a coherent group-action term.
The latter is printed as $I-\sum_{u\in G}V_u$, with the basis indices
implicit. The action convention is $L_u\ket z=\ket{uz}$ and
$R_u\ket z=\ket{zu^{-1}}$, according to the arrows.

For this permutation representation, put $S=\sum_uV_u$. Then
\begin{align}
    S^\dagger & =S,         & S^2       & =|G|S,             \\
    A         & =|G|^{-1}S, & A^\dagger & =A,    & A^2 & =A.
\end{align}
On every invariant vector, $(I-S)\psi=(1-|G|)\psi$. Except for the
trivial group, the printed expression therefore does not annihilate
the intended invariant state. The projector-complement Hamiltonian term
is $h=I-A$. This corrects the group average; it does not rescale the
color-difference tensor or assume that its state is normalized.

\section{Exact tensor and orientation}
The tensor retained from source lines 2896--2916 is
\begin{align}
    K(p,q,r,s) & =\ket{pq^{-1},qr^{-1},rs^{-1},sp^{-1}}.
\end{align}
Equation~(7.10) and its footnote, source lines 2896--2908, explicitly
prescribe clockwise rotation. We label its
virtual positions north, east, south, west by $p,q,r,s$, and its physical
positions northeast, southeast, southwest, northwest by $a,b,c,d$.
The available \path{figs6/tc-lattice.pdf} and
\path{figs6/tc-blocked-lattice.pdf} show which sites and plaquettes are
blocked and which color bonds are shared. These earlier drawings are used
only for incidence and blocking. The physical arrows and nonabelian
multiplication order used here are fixed by Equation~(7.10) and its
clockwise footnote; no literal agreement with each physical arrow of the
earlier drawings is asserted or needed. The unavailable
\path{figs6/renorm-as-diff-doubles} figure is not treated as inspected
evidence; the coordinate assignments above and all nonabelian products
are stated explicitly.

At an east--west bond the shared color is $q_L=s_R=x$. Replacing
$x$ by $ux$ acts on the actual eight physical spins as
\begin{align}
    (a,b,c,d;e,f,g,h) & \longmapsto
    (au^{-1},ub,c,d;e,f,gu^{-1},uh).
\end{align}
The hole uses the ordered spin tuple
$(NW.b,SW.a,SE.d,NE.c)$. Shared colors give these four spins as
$(xy^{-1},yz^{-1},zw^{-1},wx^{-1})$, whose product is one.
Neither identity assumes an abelian group.

\section{Local assertions and remaining obligations}
The local tensor image is the entire span of product-one computational
basis vectors. When $abcd=1$, an explicit preimage is
\begin{align}
    F(1,a^{-1},(ab)^{-1},(abc)^{-1}) & =(a,b,c,d).
\end{align}
Thus the local flatness projector is
the actual range projector of $K$, rather than an assumed tensor-image
condition.

The east--west group average fixes the actual open two-tensor contraction
before every choice of the six exterior colors: reindex its shared-color
sum by $x\mapsto ux$. The diagonal four-block plaquette projector fixes
the actual open four-tensor contraction before every choice of the eight
exterior colors, by the displayed telescoping product. Their complements
annihilate the respective contractions.

More precisely, let $\mathsf C$ be the two-block contraction map and
$P_{00}=P_0\otimes P_0$ the two local flatness constraints. Then
$P_{00}A_e=A_eP_{00}$ and
\begin{align}
    \operatorname{ran}\mathsf C & =\operatorname{ran}(A_eP_{00}).
\end{align}
For a product-one physical pair $\xi=(a,b,c,d;e,f,g,h)$, set
\begin{align}
    \beta(\xi) & =(a,b^{-1},(bc)^{-1};efg,fg,g).
\end{align}
These boundary colors reconstruct $\xi$ at shared color one.
Averaging that physical basis vector gives exactly $|G|^{-1}$ times
the corresponding contraction column. This supplies the reverse range
inclusion; tensor-image membership is not assumed.

The four-block hole projector also commutes with the coherent action
on its northern bond and with that bond's normalized average.
The hole holonomy is conjugated by $u$, so the product-one condition
is preserved even on the full ambient physical space.

Source lines 2918--2923 describe the three families of blocked Hamiltonian
terms. The present assertions are a local part of that description. The
following statements require further arguments:
\begin{enumerate}
    \item a placement of all three families on a specified full lattice,
          with every bond orientation and boundary convention accounted for;
    \item commutation of all placed terms and a full common-kernel
          identification, including the appropriate topological sectors;
    \item a comparison of that explicit common kernel with the existing
          PEPS parent Hamiltonian, rather than merely local annihilation;
    \item the actual unblocked arbitrary-group $T$-to-$K$ physical
          renormalization, with its redundant factors and normalization.
\end{enumerate}
The local physical relabelling between the review's dual tensor and the
clockwise tensor does not construct the last operation. Generic parent
classification results do not establish the second or third item unless
their local kernels are identified with the explicit Hamiltonian here.

\end{document}
